% Options for packages loaded elsewhere
% Options for packages loaded elsewhere
\PassOptionsToPackage{unicode}{hyperref}
\PassOptionsToPackage{hyphens}{url}
\PassOptionsToPackage{dvipsnames,svgnames,x11names}{xcolor}
%
\documentclass[
  letterpaper,
  DIV=11,
  numbers=noendperiod]{scrartcl}
\usepackage{xcolor}
\usepackage{amsmath,amssymb}
\setcounter{secnumdepth}{-\maxdimen} % remove section numbering
\usepackage{iftex}
\ifPDFTeX
  \usepackage[T1]{fontenc}
  \usepackage[utf8]{inputenc}
  \usepackage{textcomp} % provide euro and other symbols
\else % if luatex or xetex
  \usepackage{unicode-math} % this also loads fontspec
  \defaultfontfeatures{Scale=MatchLowercase}
  \defaultfontfeatures[\rmfamily]{Ligatures=TeX,Scale=1}
\fi
\usepackage{lmodern}
\ifPDFTeX\else
  % xetex/luatex font selection
\fi
% Use upquote if available, for straight quotes in verbatim environments
\IfFileExists{upquote.sty}{\usepackage{upquote}}{}
\IfFileExists{microtype.sty}{% use microtype if available
  \usepackage[]{microtype}
  \UseMicrotypeSet[protrusion]{basicmath} % disable protrusion for tt fonts
}{}
\makeatletter
\@ifundefined{KOMAClassName}{% if non-KOMA class
  \IfFileExists{parskip.sty}{%
    \usepackage{parskip}
  }{% else
    \setlength{\parindent}{0pt}
    \setlength{\parskip}{6pt plus 2pt minus 1pt}}
}{% if KOMA class
  \KOMAoptions{parskip=half}}
\makeatother
% Make \paragraph and \subparagraph free-standing
\makeatletter
\ifx\paragraph\undefined\else
  \let\oldparagraph\paragraph
  \renewcommand{\paragraph}{
    \@ifstar
      \xxxParagraphStar
      \xxxParagraphNoStar
  }
  \newcommand{\xxxParagraphStar}[1]{\oldparagraph*{#1}\mbox{}}
  \newcommand{\xxxParagraphNoStar}[1]{\oldparagraph{#1}\mbox{}}
\fi
\ifx\subparagraph\undefined\else
  \let\oldsubparagraph\subparagraph
  \renewcommand{\subparagraph}{
    \@ifstar
      \xxxSubParagraphStar
      \xxxSubParagraphNoStar
  }
  \newcommand{\xxxSubParagraphStar}[1]{\oldsubparagraph*{#1}\mbox{}}
  \newcommand{\xxxSubParagraphNoStar}[1]{\oldsubparagraph{#1}\mbox{}}
\fi
\makeatother


\usepackage{longtable,booktabs,array}
\usepackage{calc} % for calculating minipage widths
% Correct order of tables after \paragraph or \subparagraph
\usepackage{etoolbox}
\makeatletter
\patchcmd\longtable{\par}{\if@noskipsec\mbox{}\fi\par}{}{}
\makeatother
% Allow footnotes in longtable head/foot
\IfFileExists{footnotehyper.sty}{\usepackage{footnotehyper}}{\usepackage{footnote}}
\makesavenoteenv{longtable}
\usepackage{graphicx}
\makeatletter
\newsavebox\pandoc@box
\newcommand*\pandocbounded[1]{% scales image to fit in text height/width
  \sbox\pandoc@box{#1}%
  \Gscale@div\@tempa{\textheight}{\dimexpr\ht\pandoc@box+\dp\pandoc@box\relax}%
  \Gscale@div\@tempb{\linewidth}{\wd\pandoc@box}%
  \ifdim\@tempb\p@<\@tempa\p@\let\@tempa\@tempb\fi% select the smaller of both
  \ifdim\@tempa\p@<\p@\scalebox{\@tempa}{\usebox\pandoc@box}%
  \else\usebox{\pandoc@box}%
  \fi%
}
% Set default figure placement to htbp
\def\fps@figure{htbp}
\makeatother





\setlength{\emergencystretch}{3em} % prevent overfull lines

\providecommand{\tightlist}{%
  \setlength{\itemsep}{0pt}\setlength{\parskip}{0pt}}



 


\KOMAoption{captions}{tableheading}
\makeatletter
\@ifpackageloaded{caption}{}{\usepackage{caption}}
\AtBeginDocument{%
\ifdefined\contentsname
  \renewcommand*\contentsname{Table of contents}
\else
  \newcommand\contentsname{Table of contents}
\fi
\ifdefined\listfigurename
  \renewcommand*\listfigurename{List of Figures}
\else
  \newcommand\listfigurename{List of Figures}
\fi
\ifdefined\listtablename
  \renewcommand*\listtablename{List of Tables}
\else
  \newcommand\listtablename{List of Tables}
\fi
\ifdefined\figurename
  \renewcommand*\figurename{Figure}
\else
  \newcommand\figurename{Figure}
\fi
\ifdefined\tablename
  \renewcommand*\tablename{Table}
\else
  \newcommand\tablename{Table}
\fi
}
\@ifpackageloaded{float}{}{\usepackage{float}}
\floatstyle{ruled}
\@ifundefined{c@chapter}{\newfloat{codelisting}{h}{lop}}{\newfloat{codelisting}{h}{lop}[chapter]}
\floatname{codelisting}{Listing}
\newcommand*\listoflistings{\listof{codelisting}{List of Listings}}
\makeatother
\makeatletter
\makeatother
\makeatletter
\@ifpackageloaded{caption}{}{\usepackage{caption}}
\@ifpackageloaded{subcaption}{}{\usepackage{subcaption}}
\makeatother
\usepackage{bookmark}
\IfFileExists{xurl.sty}{\usepackage{xurl}}{} % add URL line breaks if available
\urlstyle{same}
\hypersetup{
  pdftitle={Matris İşlemlerinin Özellikleri},
  pdfauthor={Öğr. Gör. Oktay Cesur},
  colorlinks=true,
  linkcolor={blue},
  filecolor={Maroon},
  citecolor={Blue},
  urlcolor={Blue},
  pdfcreator={LaTeX via pandoc}}


\title{Matris İşlemlerinin Özellikleri}
\usepackage{etoolbox}
\makeatletter
\providecommand{\subtitle}[1]{% add subtitle to \maketitle
  \apptocmd{\@title}{\par {\large #1 \par}}{}{}
}
\makeatother
\subtitle{MATE 213 --- İlk Öğretim Bloğu}
\author{Öğr. Gör. Oktay Cesur}
\date{2026-07-24}
\begin{document}
\maketitle


\subsection{Kapsam ve Kabuller}\label{kapsam-ve-kabuller}

\[
A=[a_{ij}],\quad B=[b_{ij}],\quad C=[c_{ij}],
\qquad
\alpha,\beta\in\mathbb{R}
\]

\begin{itemize}
\tightlist
\item
  Bütün eşitlikler boyut uyumu varsayımı altında
\item
  \(0\): uygun şekilli sıfır matrisi
\item
  \(I_n\): \(n\times n\) birim matris
\end{itemize}

Bu not, önceki notlarda tanımlanan dört işlemin cebirsel kurallarını bir
arada topluyor: toplama, skalerle çarpma, matris--vektör çarpımı ve
matris--matris çarpımı. Her kural önce genel biçimde yazılıyor, sonra
küçük matrislerle sayısal olarak doğrulanıyor. Geçersiz kurallar için
karşı örnek veriliyor.

Aşağıdaki bütün eşitliklerde boyutların işlemi tanımlı kılacak biçimde
seçildiği varsayılıyor. Örneğin \(A(B+C)=AB+AC\) eşitliğinde \(B\) ile
\(C\) aynı şekilde olmalı, \(A\) matrisinin sütun sayısı da bu şeklin
satır sayısına eşit olmalıdır. \(0\) ve \(I\) sembolleri de her
kullanımda ilgili işlemi tanımlı kılan şekilde okunur; \(I_mA=AI_n=A\)
eşitliğinde iki birim matrisin boyutu farklıdır.

Matris--vektör çarpımı ayrı bir başlık altında ele alınmıyor. Bir sütun
vektörü \(n\times1\) şekilli matris olduğundan, aşağıdaki kuralların
hepsi \(p=1\) durumunda doğrudan \(Ax\) ifadesine uygulanır.

\begin{center}\rule{0.5\linewidth}{0.5pt}\end{center}

\subsection{İşlemlerin Özeti}\label{iux15flemlerin-uxf6zeti}

\begin{longtable}[]{@{}lll@{}}
\toprule\noalign{}
İşlem & Eşleşen yapı & Sonuç \\
\midrule\noalign{}
\endhead
\bottomrule\noalign{}
\endlastfoot
\(A+B\) & aynı konumlar & aynı şekilli matris \\
\(\alpha A\) & tek skaler--bütün elemanlar & aynı şekilli matris \\
\(Ax\) & satırlar--tek katsayı sütunu & tek çıktı sütunu \\
\(AB\) & satırlar--çoklu katsayı sütunları & çoklu çıktı matrisi \\
\end{longtable}

Dört işlemin tanımlılık koşulu aynı soruyla bulunur: hangi bileşenler ya
da eksenler eşleşiyor? Toplamada bütün konumlar bire bir eşleşir, bu
nedenle şekiller tamamen aynı olmalıdır. Skalerle çarpmada eşleşme
koşulu yoktur; tek sayı bütün elemanlara uygulanır. Çarpımda ise yalnız
\(A\) matrisinin sütun ekseni ile \(B\) matrisinin satır ekseni eşleşir.

Koşulların farklı olmasının doğrudan bir sonucu var: \(2\times3\) ve
\(3\times2\) şekilli iki matris toplanamaz, fakat uygun sırada
çarpılabilir. Tersine, aynı şekilli iki dikdörtgen matris toplanabilir
ama çarpılamaz.

Bu ayrım, aşağıdaki özellik listelerinin neden ayrı ayrı yazıldığını da
açıklıyor. Toplamanın özellikleri eleman bazlı işlemlerden, çarpımın
özellikleri ise çarp-topla yapısından geliyor.

\begin{center}\rule{0.5\linewidth}{0.5pt}\end{center}

\subsection{Toplama ve Skalerle
Çarpma}\label{toplama-ve-skalerle-uxe7arpma}

\[
A+B=B+A,
\qquad
(A+B)+C=A+(B+C)
\]

\[
A+0=A,
\qquad
A+(-A)=0
\]

\[
\alpha(A+B)=\alpha A+\alpha B,
\qquad
(\alpha+\beta)A=\alpha A+\beta A,
\qquad
\alpha(\beta A)=(\alpha\beta)A
\]

Bu eşitliklerin tümü eleman düzeyinde doğrudan doğrulanır. Örneğin
değişme özelliği için
\((A+B)_{ij}=a_{ij}+b_{ij}=b_{ij}+a_{ij}=(B+A)_{ij}\) yazmak yeterlidir;
matrisler için geçerli olan şey, gerçek sayılarda toplamanın değişme
özelliğinden geliyor. Aynı gerekçe birleşme ve dağılma eşitlikleri için
de geçerlidir.

\(0\) matrisi toplamanın etkisiz elemanı, \(-A\) ise \(A\) matrisinin
toplamsal tersidir. Bu iki nesnenin varlığı, \(\mathbb{R}^{m\times n}\)
kümesinin toplama ve skalerle çarpma altında kapalı bir yapı
oluşturmasını sağlar. Vektör uzayı tanımını gördüğümüzde bu listenin tam
olarak aksiyomlarla örtüştüğünü fark edeceksiniz.

Çarpımda durum farklıdır. Aşağıdaki listede değişme özelliği yer almıyor
ve sadeleştirme kuralı da geçerli değil. Farkın kaynağı, çarpımın eleman
bazlı bir işlem olmamasıdır.

\begin{center}\rule{0.5\linewidth}{0.5pt}\end{center}

\subsection{Çarpımın Cebirsel
Özellikleri}\label{uxe7arpux131mux131n-cebirsel-uxf6zellikleri}

\[
(AB)C=A(BC)
\]

\[
A(B+C)=AB+AC,
\qquad
(A+B)C=AC+BC
\]

\[
\alpha(AB)=(\alpha A)B=A(\alpha B)
\]

\[
A0=0,
\qquad
0A=0
\]

Bu liste, matris çarpımının gerçek sayılardaki çarpımdan hangi yönlerden
ayrıldığını belirler. Birleşme, dağılma ve skalerle uyum korunur.
Korunmayan iki temel özellik değişme ve sadeleştirmedir; ikisini de
ayrıca ele alacağız.

Dağılma iki ayrı eşitlik olarak yazılıyor. \(A(B+C)\) ile \((B+C)A\)
farklı ifadelerdir ve soldan dağıtmakla sağdan dağıtmak birbirinin
yerine kullanılamaz. Sayılarda tek bir dağılma kuralı yazmak yeterliyken
burada iki kural gerekiyor, çünkü çarpanların sırası serbest değil.

Sıfır matrisiyle çarpım kuralı da göründüğünden dardır. \(A0=0\)
eşitliği doğrudur, fakat tersi yönde bir çıkarım yapılamaz: \(AB=0\)
olması \(A=0\) ya da \(B=0\) olmasını gerektirmez. Bu duruma birazdan
sayısal bir karşı örnekle geleceğiz.

\begin{center}\rule{0.5\linewidth}{0.5pt}\end{center}

\subsection{Birleşme Özelliğinin
Doğrulaması}\label{birleux15fme-uxf6zelliux11finin-doux11frulamasux131}

\[
\bigl((AB)C\bigr)_{ij}
=\sum_{k}\Bigl(\sum_{\ell}a_{i\ell}b_{\ell k}\Bigr)c_{kj}
\]

\[
=\sum_{\ell}a_{i\ell}\Bigl(\sum_{k}b_{\ell k}c_{kj}\Bigr)
=\bigl(A(BC)\bigr)_{ij}
\]

Birleşme özelliği tek satırlık bir kabul değil, indis hesabıyla
gösterilebilen bir sonuçtur. Sol taraftaki \((i,j)\) elemanı yazılırken
önce \(AB\) çarpımının \((i,k)\) elemanı
\(\sum_\ell a_{i\ell}b_{\ell k}\) olarak açılır, sonra \(C\) matrisiyle
çarpım için \(k\) üzerinden toplanır.

İkinci adımda iki sonlu toplamın sırası değiştirilir. Bu adım geçerlidir
çünkü toplanan terim sayısı sonludur ve her terim
\(a_{i\ell}b_{\ell k}c_{kj}\) biçiminde üç skalerin çarpımıdır. Sıra
değiştikten sonra parantez içinde kalan ifade \((BC)_{\ell j}\)
elemanıdır; dolayısıyla sağ taraf \(A(BC)\) çarpımının \((i,j)\)
elemanına eşittir.

Bütün \((i,j)\) çiftleri için eşitlik sağlandığından iki matris eşittir.
Aynı yöntem dağılma eşitlikleri için de kullanılır; orada toplamı
ayırmak yeterli olduğundan hesap daha kısadır.

\begin{center}\rule{0.5\linewidth}{0.5pt}\end{center}

\subsection{Sayısal Doğrulama:
Birleşme}\label{sayux131sal-doux11frulama-birleux15fme}

\[
A=\begin{bmatrix}1&2\\0&3\end{bmatrix},
\quad
B=\begin{bmatrix}2&1\\1&0\end{bmatrix},
\quad
C=\begin{bmatrix}1&0\\2&1\end{bmatrix}
\]

\[
AB=\begin{bmatrix}4&1\\3&0\end{bmatrix}
\ \Longrightarrow\
(AB)C=\begin{bmatrix}6&1\\3&0\end{bmatrix}
\]

\[
BC=\begin{bmatrix}4&1\\1&0\end{bmatrix}
\ \Longrightarrow\
A(BC)=\begin{bmatrix}6&1\\3&0\end{bmatrix}
\]

Soldan hesapta önce \(AB\) bulunur: birinci satır \((1,2)\) ile birinci
sütun \((2,1)\) çarpılıp toplandığında \(1(2)+2(1)=4\), ikinci sütun
\((1,0)\) ile \(1(1)+2(0)=1\) elde edilir. İkinci satır \((0,3)\) için
değerler \(3\) ve \(0\) çıkar. Bulunan \(AB\) matrisi \(C\) ile
çarpıldığında \((1,1)\) elemanı \(4(1)+1(2)=6\) olur.

Sağdan hesapta önce \(BC\) bulunur ve \(A\) ile çarpılır: \((1,1)\)
elemanı \(1(4)+2(1)=6\) çıkar. İki yol da aynı matrisi veriyor, ama ara
sonuçlar farklıdır. \(AB\) ile \(BC\) matrisleri birbirine eşit değil;
eşit olan yalnız üçlü çarpımın sonucudur.

Ara sonuçların farklı olması hesap yükü açısından işe yarar. Boyutlar
eşit olmadığında iki yoldan biri belirgin biçimde daha az çarpma
gerektirebilir; örneğin \(A\) matrisi \(100\times2\), \(B\) matrisi
\(2\times100\) ve \(C\) matrisi \(100\times1\) ise \(A(BC)\) hesabı
\((AB)C\) hesabından çok daha kısadır.

\begin{center}\rule{0.5\linewidth}{0.5pt}\end{center}

\subsection{Sayısal Doğrulama:
Dağılma}\label{sayux131sal-doux11frulama-daux11fux131lma}

\[
B+C=\begin{bmatrix}3&1\\3&1\end{bmatrix}
\ \Longrightarrow\
A(B+C)=\begin{bmatrix}9&3\\9&3\end{bmatrix}
\]

\[
AB=\begin{bmatrix}4&1\\3&0\end{bmatrix},
\qquad
AC=\begin{bmatrix}5&2\\6&3\end{bmatrix}
\]

\[
AB+AC=\begin{bmatrix}9&3\\9&3\end{bmatrix}
\]

Aynı üç matrisle dağılma eşitliğini doğruluyoruz. Soldan hesapta önce
\(B+C\) toplamı alınır; toplama eleman bazlı olduğu için \((1,1)\)
elemanı \(2+1=3\) olur. Bulunan matris \(A\) ile çarpıldığında \((1,1)\)
elemanı \(1(3)+2(3)=9\) çıkar.

Sağdan hesapta iki çarpım ayrı ayrı yapılır. \(AC\) matrisinin \((1,1)\)
elemanı \(1(1)+2(2)=5\), \(AB\) matrisininki \(4\)'tür; toplamları \(9\)
eder. Bütün konumlarda aynı eşleşme sağlandığı için iki taraf eşittir.

Bu eşitliğin sağdan biçimi ayrıca kontrol edilmelidir. \((B+C)A\) hesabı
\(A(B+C)\) hesabına eşit değildir; burada \(B+C\) ile \(A\)
matrislerinin çarpım sırası değiştiğinden sonuç da değişir. Dağılma
kuralı sırayı serbest bırakmaz, yalnız parantezi açar.

\begin{center}\rule{0.5\linewidth}{0.5pt}\end{center}

\subsection{Birim ve Sıfır
Matrisi}\label{birim-ve-sux131fux131r-matrisi}

\[
I_mA=AI_n=A
\qquad (A:m\times n)
\]

\[
\begin{bmatrix}1&0\\0&1\end{bmatrix}
\begin{bmatrix}1&2\\0&3\end{bmatrix}
=\begin{bmatrix}1&2\\0&3\end{bmatrix}
\]

\[
A0=0,
\qquad
0A=0
\]

Birim matris çarpımın etkisiz elemanıdır, fakat toplamadaki \(0\)
matrisinden farklı olarak sağdan ve soldan farklı boyutlarda kullanılır.
\(A\) matrisi \(m\times n\) şekilliyse soldan çarpan \(I_m\), sağdan
çarpan \(I_n\) olmalıdır. Kare olmayan bir matris için bu iki birim
matris aynı nesne değildir.

\(I_mA=A\) eşitliği tanımdan çıkar. \(I_m\) matrisinin \(i\). satırında
yalnız \(i\). konumda \(1\), diğer konumlarda \(0\) vardır; çarp-topla
hesabında bu satır \(A\) matrisinin yalnız \(i\). satırını seçer, kalan
terimler sıfırla çarpıldığı için düşer.

Sıfır matrisi çarpımı yutar. Ancak bu kuralın tersi geçerli değildir;
\(AB=0\) eşitliğinden çarpanlar hakkında sonuç çıkarılamaz. Sıradaki
karşı örnek bunu gösteriyor.

\begin{center}\rule{0.5\linewidth}{0.5pt}\end{center}

\subsection{\texorpdfstring{Sıfır Bölen: \(AB=0\) Ne Demek
Değildir?}{Sıfır Bölen: AB=0 Ne Demek Değildir?}}\label{sux131fux131r-buxf6len-ab0-ne-demek-deux11fildir}

\[
A=\begin{bmatrix}1&1\\1&1\end{bmatrix},
\qquad
B=\begin{bmatrix}1&1\\-1&-1\end{bmatrix}
\]

\[
AB=\begin{bmatrix}0&0\\0&0\end{bmatrix},
\qquad
A\ne0,\quad B\ne0
\]

Çarpımın \((1,1)\) elemanı \(1(1)+1(-1)=0\), \((1,2)\) elemanı
\(1(1)+1(-1)=0\) olur; ikinci satır da aynı değerleri verir çünkü \(A\)
matrisinin iki satırı aynıdır. Sonuç sıfır matrisidir, oysa çarpanların
hiçbiri sıfır değildir.

Gerçek sayılarda \(ab=0\) eşitliği çarpanlardan birinin sıfır olmasını
gerektirir; bu, denklem çözerken sürekli kullandığımız bir çıkarımdır.
Matrislerde aynı çıkarımı yapmak hatalıdır. Örneğin \(A^2=A\)
denkleminden \(A(A-I)=0\) elde edilir, fakat buradan \(A=0\) ya da
\(A=I\) sonucu çıkmaz.

Karşı örnekteki iki matrisin ortak yanı sütunlarının birbirinin katı
olmasıdır. Bu durumun ne anlama geldiğini lineer bağımsızlık ve rank
konularında göreceğiz; şimdilik kuralın geçersiz olduğunu bilmek
yeterli.

\begin{center}\rule{0.5\linewidth}{0.5pt}\end{center}

\subsection{Sadeleştirme Kuralı
Geçersizdir}\label{sadeleux15ftirme-kuralux131-geuxe7ersizdir}

\[
A=\begin{bmatrix}1&1\\1&1\end{bmatrix},
\qquad
B=\begin{bmatrix}1&0\\0&1\end{bmatrix},
\qquad
C=\begin{bmatrix}0&1\\1&0\end{bmatrix}
\]

\[
AB=AC=\begin{bmatrix}1&1\\1&1\end{bmatrix},
\qquad
B\ne C
\]

\(B\) birim matris olduğu için \(AB=A\) olur. \(AC\) hesabında ise \(C\)
matrisi \(A\) matrisinin sütunlarının yerini değiştirir; \(A\)
matrisinin iki sütunu aynı olduğundan sonuç yine \(A\) çıkar. İki çarpım
eşit, fakat \(B\) ile \(C\) farklı matrislerdir.

Sayılarda \(ab=ac\) ve \(a\ne0\) ise \(b=c\) sonucunu çıkarırız.
Matrislerde \(A\ne0\) koşulu bu çıkarım için yeterli değildir. Eşitliğin
iki tarafını \(A\) ile ``sadeleştirmek'', ancak \(A\) matrisinin tersi
varsa yapılabilir; tersinirlik konusunda bu koşulu tam olarak
belirleyeceğiz.

İki karşı örnekteki \(A\) matrisi aynıdır ve bu tesadüf değildir. Sıfır
bölen olma ile sadeleştirilememe aynı olgunun iki yüzüdür: \(AB=AC\)
eşitliği \(A(B-C)=0\) demektir, yani \(A\) sıfır olmayan bir matrisi
sıfıra götürüyor.

\begin{center}\rule{0.5\linewidth}{0.5pt}\end{center}

\subsection{Transpoz ve Çarpım}\label{transpoz-ve-uxe7arpux131m}

\[
(A+B)^T=A^T+B^T,
\qquad
(\alpha A)^T=\alpha A^T,
\qquad
(A^T)^T=A
\]

\[
\boxed{(AB)^T=B^TA^T}
\]

\begin{quote}
Transpoz alındığında çarpanların sırası ters döner.
\end{quote}

İlk üç eşitlik eleman bazlı işlemlerle uyumludur ve doğrudan tanımdan
çıkar. Dördüncü eşitlik ise sırayı ters çevirir. Nedeni boyutlarda
görünür: \(A\) matrisi \(m\times n\), \(B\) matrisi \(n\times p\) ise
\(AB\) çarpımı \(m\times p\), transpozu \(p\times m\) şekillidir. Sağ
tarafta \(B^T\) matrisi \(p\times n\), \(A^T\) matrisi \(n\times m\)
olduğundan çarpım \(p\times m\) çıkar ve şekiller uyar.

Sırayı korumaya çalışmak çoğu zaman tanımsız bir ifade üretir.
\(A^TB^T\) çarpımı \((n\times m)(p\times n)\) biçiminde olur ve \(m=p\)
değilse hesaplanamaz bile. Hesaplanabildiği durumlarda da genel olarak
\((AB)^T\) matrisine eşit değildir.

Eleman düzeyinde gerekçe kısadır:
\(\bigl((AB)^T\bigr)_{ij}=(AB)_{ji}=\sum_k a_{jk}b_{ki}\) olur. Sağ
tarafta ise
\((B^TA^T)_{ij}=\sum_k (B^T)_{ik}(A^T)_{kj}=\sum_k b_{ki}a_{jk}\) elde
edilir; iki toplam aynı terimlerden oluşur.

\begin{center}\rule{0.5\linewidth}{0.5pt}\end{center}

\subsection{Sayısal Doğrulama:
Transpoz}\label{sayux131sal-doux11frulama-transpoz}

\[
AB=\begin{bmatrix}4&1\\3&0\end{bmatrix}
\ \Longrightarrow\
(AB)^T=\begin{bmatrix}4&3\\1&0\end{bmatrix}
\]

\[
B^TA^T=
\begin{bmatrix}2&1\\1&0\end{bmatrix}
\begin{bmatrix}1&0\\2&3\end{bmatrix}
=\begin{bmatrix}4&3\\1&0\end{bmatrix}
\]

\[
A^TB^T=\begin{bmatrix}2&1\\7&2\end{bmatrix}
\]

Aynı \(A\) ve \(B\) matrisleriyle devam ediyoruz. \(B^T\) matrisi \(B\)
simetrik olduğu için kendisine eşittir; \(A^T\) matrisi ise
\(\begin{bmatrix}1&0\\2&3\end{bmatrix}\) olur. Çarpımın \((1,1)\)
elemanı \(2(1)+1(2)=4\), \((1,2)\) elemanı \(2(0)+1(3)=3\) çıkar ve
sonuç \((AB)^T\) matrisine eşittir.

Ters sırada yazıldığında sonuç değişir. \(A^TB^T\) çarpımının \((2,1)\)
elemanı \(2(2)+3(1)=7\) olur, oysa \((AB)^T\) matrisinin aynı
konumundaki değer \(1\)'dir. Burada iki matris de \(2\times2\) olduğu
için çarpım tanımlıdır; eşit olmamaları boyut sorunundan değil, sıranın
kendisinden kaynaklanıyor.

Kare olmayan matrislerde hata daha erken yakalanır. \(A\) matrisi
\(2\times3\), \(B\) matrisi \(3\times2\) ise \(A^TB^T\) çarpımı
\((3\times2)(2\times3)\) biçiminde olur ve \(3\times3\) bir matris
verir; \((AB)^T\) ise \(2\times2\) şekillidir. Boyut kontrolü bu durumda
yanlış sırayı doğrudan gösterir.

\begin{center}\rule{0.5\linewidth}{0.5pt}\end{center}

\subsection{Matrisin Kuvvetleri}\label{matrisin-kuvvetleri}

\[
A^0=I,
\qquad
A^k=\underbrace{AA\cdots A}_{k\ \text{çarpan}}
\qquad (A\ \text{kare})
\]

\[
A=\begin{bmatrix}1&1\\0&1\end{bmatrix}
\ \Longrightarrow\
A^2=\begin{bmatrix}1&2\\0&1\end{bmatrix},
\quad
A^3=\begin{bmatrix}1&3\\0&1\end{bmatrix}
\]

\[
N=\begin{bmatrix}0&1\\0&0\end{bmatrix}
\ \Longrightarrow\
N^2=0
\]

Kuvvet yalnız kare matrisler için tanımlıdır; \(A\) matrisi kare değilse
\(AA\) çarpımının iç boyutları uyuşmaz. Birleşme özelliği sayesinde
çarpanların hangi sırayla gruplandığı önemsizdir, bu yüzden
\(A^{k+\ell}=A^kA^\ell\) eşitliği geçerlidir.

İlk örnekte üst köşedeki eleman her adımda \(1\) artıyor; tümevarımla
\(A^k=\begin{bmatrix}1&k\\0&1\end{bmatrix}\) olduğu gösterilebilir.
İkinci örnekte ise sıfırdan farklı bir matrisin karesi sıfır çıkıyor.
Sayılarda \(x^2=0\) denklemi yalnız \(x=0\) çözümünü verirken
matrislerde durum farklıdır; bu, sıfır bölen örneğinin kuvvet
biçimindeki karşılığıdır.

Değişme özelliğinin olmaması kuvvet hesaplarında da görünür.
\((AB)^2=ABAB\) ifadesi ancak \(AB=BA\) olduğunda \(A^2B^2\) biçiminde
toplanabilir. Aynı nedenle \((A+B)^2=A^2+AB+BA+B^2\) olur; orta terimler
genel olarak \(2AB\) vermez.

\begin{center}\rule{0.5\linewidth}{0.5pt}\end{center}

\subsection{Sık Yapılan Hatalar}\label{sux131k-yapux131lan-hatalar}

\begin{enumerate}
\def\labelenumi{\arabic{enumi}.}
\tightlist
\item
  \(AB=0\) eşitliğinden bir çarpanı sıfır saymak.
\item
  \(AB=AC\) eşitliğinde \(A\) ile sadeleştirmek.
\item
  \((AB)^T\) yerine \(A^TB^T\) yazmak.
\item
  \((A+B)^2\) açılımında \(2AB\) yazmak.
\item
  Dağılmayı sıra değiştirmek için kullanmak.
\end{enumerate}

Birinci ve ikinci hatalar aynı kaynaktan gelir: gerçek sayılardaki
çarpımın tersinir olmasına dayanan çıkarımlar matrislere taşınıyor.
İkisi de ancak ilgili matrisin tersi varsa geçerlidir ve bu koşul genel
olarak sağlanmaz.

Üçüncü hata sıra değişimini atlar. Boyutlar farklıysa yanlış ifade zaten
tanımsız çıkar ve hata hemen görünür; kare matrislerde ise hesap
tamamlanır ve yanlış sonuç fark edilmeyebilir. Dördüncü hata da
benzerdir: \(AB\) ile \(BA\) genel olarak eşit olmadığı için orta
terimler birleştirilemez.

Beşinci hata dağılma kuralının ne söylediğini genişletir.
\(A(B+C)=AB+AC\) eşitliği parantezi açar, çarpanların sırasını serbest
bırakmaz. \((B+C)A\) ifadesi açıldığında \(BA+CA\) elde edilir ve bu
genel olarak \(AB+AC\) toplamına eşit değildir.

\begin{center}\rule{0.5\linewidth}{0.5pt}\end{center}

\subsection{Karar Soruları}\label{karar-sorularux131}

Hesap yapmadan karar verin:

\begin{enumerate}
\def\labelenumi{\arabic{enumi}.}
\tightlist
\item
  \(A:3\times4\) için \(A^2\) tanımlı mı?
\item
  \(AB=BA\) her zaman yanlış mı?
\item
  \(A\ne0\), \(AB=0\) ise \(B=0\) mı?
\item
  \((ABC)^T\) nasıl yazılır?
\end{enumerate}

Birinci soruda \(A^2=AA\) çarpımı \((3\times4)(3\times4)\) biçiminde
olur; iç boyutlar \(4\) ve \(3\) uyuşmadığı için tanımsızdır. Kuvvet
yalnız kare matrislerde tanımlıdır.

İkinci soruda cevap hayırdır. Değişme özelliği genel olarak geçerli
değildir, fakat özel çiftlerde sağlanabilir: \(A\) ile \(I\), \(A\) ile
\(A^k\), ya da iki köşegen matris her zaman yer değiştirebilir. ``Genel
olarak geçerli değil'' ile ``hiçbir zaman geçerli değil'' farklı
ifadelerdir.

Üçüncü soruda cevap yine hayırdır; karşı örneği bu notta gördük.
Dördüncü soruda kural iki kez uygulanır:
\((ABC)^T=\bigl((AB)C\bigr)^T=C^T(AB)^T=C^TB^TA^T\). Sıra baştan sona
tersine döner.

\begin{center}\rule{0.5\linewidth}{0.5pt}\end{center}

\subsection{Sonraki Adım: Kural Olarak
Matris}\label{sonraki-adux131m-kural-olarak-matris}

Elimizde:

\[
\text{tanım}+\text{mekanizma}+\text{cebirsel kurallar}
\]

\begin{itemize}
\tightlist
\item
  Açık kalan: \(x\mapsto Ax\) neyi temsil ediyor?
\item
  Sonraki not: dönüşüm okuması
\end{itemize}

Bu notta matris işlemlerinin cebirsel yapısını tamamladık. Toplama ve
skalerle çarpma sayılardaki karşılıklarıyla aynı kuralları izliyor;
çarpım ise birleşme, dağılma ve skalerle uyumu koruyup değişme ile
sadeleştirmeyi kaybediyor. Kaybedilen iki kuralın karşı örneklerini
sayısal olarak gördük.

Bundan sonra \(A\) matrisini bir hesap düzeni olarak değil, \(x\)
vektörünü \(Ax\) vektörüne götüren bir kural olarak okuyacağız. Buradaki
cebirsel özellikler o okumada yeniden karşımıza çıkacak:
\(A(u+v)=Au+Av\) eşitliği kuralın toplamayı koruduğunu, \((AB)x=A(Bx)\)
eşitliği ise çarpımın ardışık uygulamaya karşılık geldiğini söyleyecek.




\end{document}
